Neste estudo, foi desenvolvida uma metodologia para a classificação automática dos graus ISUP em imagens histopatológicas de próstata, considerando desafios relevantes como o desbalanceamento de classes, o ruído nos rótulos e a natureza ordinal da progressão tumoral. A introdução de uma função de perda híbrida, combinando regressão ordinal e perda focal, contribuiu para melhorar o desempenho do modelo ao capturar explicitamente a hierarquia entre os graus de malignidade. Destaca-se que o maior ganho de desempenho foi obtido com a aplicação da perda ordinal, que aplicada em conjunto com a remoção de ruído elevou o índice $\kappa_{\text{quad}}$ em aproximadamente 2,5\%, mantendo elevada estabilidade (desvio padrão de 0,009). De forma geral, a abordagem proposta elevou o índice $\kappa_{\text{quad}}$ de 0,826 para 0,857 (ganho de aproximadamente 3\%) e a acurácia de 0,592 para 0,669 (ganho de cerca de 7\%) em relação ao \textit{baseline}. Na comparação entre arquiteturas da família EfficientNet, observou-se que o aumento da complexidade do modelo não resultou em melhor desempenho, havendo inclusive leve degradação nas variantes mais profundas. Esses resultados indicam que arquiteturas mais complexas requerem ajustes mais refinados de treinamento para alcançar seu potencial máximo. Apesar dos excelentes resultados, a principal limitação do estudo reside no uso de um único banco de dados, ainda que multicêntrico e heterogêneo, o que indica a necessidade de validações futuras em bases de dados externas para comprovar a robustez e a generalização do método.


Como trabalhos futuros, propõe-se investigar arquiteturas mais recentes baseadas em \textit{Transformers}, especialmente \textit{Vision Transformers} (ViT) e modelos híbridos \textit{CNN-Transformer}, que podem capturar melhor dependências globais em imagens histopatológicas bem como técnicas de \textit{ensemble learning} para aumentar a robustez e a capacidade de generalização do modelo. Essas abordagens têm potencial para melhorar o desempenho do sistema, sobretudo na classificação dos graus intermediários do ISUP, contribuindo para ampliar sua aplicabilidade como ferramenta de apoio à decisão diagnóstica em diferentes contextos clínicos.
